% Propietario conservado: /Users/ruben/Documents/New project/output/EXTREMA_MEDIA_RAZON_RECIPROCIDAD_ES_EN_20260918/ARTICULO/ES/source/sections/01_hilbert_nonadico.tex
% Procedencia: integral 2249, cap. 23, pp. 601--608; propietario base_c21_sin_encabezado.tex.
% Recuperación y desarrollo de pruebas. Requiere el núcleo común APP--TRIT--TPK.
\section{Realización hilbertiana de la fase nonádica y conservación de la unidad}
\label{x_reciprocidad:iv:hilbert}

La construcción de Holografía Modular Triádica proporciona primero los
estados, sus transiciones y los mapas de restricción. La realización
hilbertiana añade a sus coordenadas discretas una estructura lineal y un
producto interno. Este orden determina qué representa cada operador:
las etiquetas de la base proceden de la construcción nonádica; las
operaciones lineales actúan después sobre esas etiquetas. La memoria
genealógica permanece en el estado del TPK y acompaña a las proyecciones
que se utilicen para representarlo.

\subsection{Fibra de fase, base etiquetada y producto interno}

Para una profundidad entera \(d\geq1\), sea
\begin{equation}
 \mathcal C_{9,d}=(\ZZ/9\ZZ)^d,
 \qquad \mathcal H_d=\ell^2(\mathcal C_{9,d}),
 \qquad \mathcal A_d=B(\mathcal H_d).
 \label{x_reciprocidad:iv:eq:hilbert-finito}
\end{equation}
El producto interno es
\(\langle f,g\rangle=\sum_x\overline{f(x)}g(x)\). Las funciones
\(\delta_x\) forman una base ortonormal etiquetada por las coordenadas
de fase. Así, \(\dim\mathcal H_d=9^d\) y
\(\mathcal A_d\cong M_{9^d}(\CC)\). El álgebra matricial representa
esta fibra finita; las rutas, los cocientes y la memoria que una
proyección de fase omita siguen perteneciendo al estado de partida.

La factorización de coordenadas induce el unitario
\begin{equation}
 \mathcal H_{d+1}\longrightarrow\mathcal H_d\otimes\CC^9,
 \qquad \delta_{(x,r)}\longmapsto\delta_x\otimes\delta_r.
 \label{x_reciprocidad:iv:eq:factorizacion-hilbert}
\end{equation}
Definimos la traza normalizada
\(\tau_d(a)=9^{-d}\operatorname{Tr}(a)\). La cuenta finita y la
normalización se distinguen: \(\operatorname{Tr}(I)=9^d\), mientras
que \(\tau_d(I)=1\). Asimismo, el grupo \((\ZZ/9\ZZ)^3\) y el
espacio vectorial \(\FF_3^6\) tienen ambos \(729\) elementos, pero
conservan leyes algebraicas diferentes. La realización de una fibra de
fase no sustituye la codificación ternaria de las emisiones APP.

\subsection{Traslación y fase: representación finita de Weyl}

En la realización compleja posterior a las construcciones de
\(\pi\) y de la estructura \(\mathrm i^2=-1\), fijamos
\(\omega=\exp(2\pi\mathrm i/9)\). Sobre \(\mathcal H_1\), sean
\begin{equation}
 (Xf)(r)=f(r-1),\qquad (Zf)(r)=\omega^r f(r).
 \label{x_reciprocidad:iv:eq:weyl-operadores}
\end{equation}
La primera operación desplaza la coordenada de fase; la segunda la
evalúa mediante un carácter. Ambas son unitarias y satisfacen
\begin{equation}
 X^9=Z^9=I,\qquad ZX=\omega XZ.
 \label{x_reciprocidad:iv:eq:weyl-relacion}
\end{equation}
En efecto, \((ZXf)(r)=\omega^r f(r-1)\), mientras que
\((XZf)(r)=\omega^{r-1}f(r-1)\). El carácter aparece aquí como
realización de la fase ya construida y no selecciona los registros
numéricos del TPK.

% RESULTADO_RECUPERADO: 14_numero_heisenberg_y_d108.tex:741--788.
% Prueba algebraica desarrollada con la convención X^a Z^b ya fijada aquí.
\subsubsection{Área orientada y extensión central de las traslaciones}
\label{x_reciprocidad:iv:weyl-area-extension}

El soporte de traslaciones es \(X_9=(\ZZ/9\ZZ)^2\).
La curvatura mixta de APP construida en la
sección~\ref{x_reciprocidad:iv:extension-curvatura-mixta} asigna \(+1\) a la
plaqueta orientada de la pareja suma--producto y \(-1\) a la
pareja de orden contrario. Su extensión bilineal a dos desplazamientos
\(u=(a,b)\) y \(v=(c,d)\) es la forma alternante
\begin{equation}
 \sigma_{\mathrm{ar}}(u,v)=ad-bc\quad\text{en }\ZZ/9\ZZ.
 \label{x_reciprocidad:iv:eq:weyl-area-orientada}
\end{equation}
En efecto, la alternancia anula la evaluación sobre dos desplazamientos
del mismo eje. Los valores sobre la base orientada son
\(\sigma_{\mathrm{ar}}((1,0),(0,1))=1\) y
\(\sigma_{\mathrm{ar}}((0,1),(1,0))=-1\); distribuir ambos
argumentos da \eqref{x_reciprocidad:iv:eq:weyl-area-orientada}. La misma fórmula es
la suma de plaquetas de cualquier cadena celular orientada que realice
ese paralelogramo: las aristas interiores se cancelan. Modificar un
representante entero por un múltiplo de nueve conserva su publicación
en \(\ZZ/9\ZZ\).

El orden \(X^aZ^b\) de los operadores ya definidos fija una sección
de la extensión central. Denotemos su factor de composición por
\begin{equation}
 c_{\mathrm W}(u,v)=bc\quad\text{en }\ZZ/9\ZZ.
 \label{x_reciprocidad:iv:eq:weyl-cociclo-polarizado}
\end{equation}
Este cociclo conserva el orden de las dos operaciones, mientras
\(\sigma_{\mathrm{ar}}\) conserva la orientación del paralelogramo.

\begin{proposition}[Extensión central y signo del conmutador]
\label{x_reciprocidad:iv:prop:heisenberg-orientado}
El conjunto \(\mathfrak H_9=(\ZZ/9\ZZ)^3\), con producto
\begin{equation}
 (a,b,t)(c,d,s)=(a+c,b+d,t+s+bc),
 \label{x_reciprocidad:iv:eq:heisenberg-ley}
\end{equation}
es un grupo de orden \(729\). La aplicación
\[
 \varrho(a,b,t)=\omega^tX^aZ^b
\]
es una representación unitaria fiel sobre \(\mathcal H_1\).
Con el convenio \([g,h]=ghg^{-1}h^{-1}\), se tiene
\begin{equation}
 [\varrho(a,b,0),\varrho(c,d,0)]
 =\omega^{bc-ad}I
 =\omega^{-\sigma_{\mathrm{ar}}((a,b),(c,d))}I.
 \label{x_reciprocidad:iv:eq:weyl-conmutador-orientado}
\end{equation}
El centro y el grupo de traslaciones forman la sucesión exacta
\[
 0\longrightarrow\ZZ/9\ZZ
 \xrightarrow{\ t\mapsto(0,0,t)\ }\mathfrak H_9
 \xrightarrow{\ (a,b,t)\mapsto(a,b)\ }X_9
 \longrightarrow0.
\]
\end{proposition}
\begin{proof}
Para \(w=(e,f)\), la identidad de cociclo se escribe
\[
 c_{\mathrm W}(u,v)+c_{\mathrm W}(u+v,w)
 =bc+(b+d)e
 =de+b(c+e)
 =c_{\mathrm W}(v,w)+c_{\mathrm W}(u,v+w).
\]
Prueba la asociatividad de \eqref{x_reciprocidad:iv:eq:heisenberg-ley}; la identidad
es \((0,0,0)\) y la inversa de \((a,b,t)\) es
\((-a,-b,-t+ab)\). Los tres residuos son libres, de donde el orden
\(9^3=729\). De \(ZX=\omega XZ\) resulta
\[
 (X^aZ^b)(X^cZ^d)=\omega^{bc}X^{a+c}Z^{b+d},
\]
que demuestra la propiedad de representación. Si
\(\omega^tX^aZ^b=I\), su acción sobre \(\delta_r\) obliga primero
a \(a=0\). La evaluación en \(r=0\) da \(t=0\), y la
evaluación en \(r=1\) da \(b=0\), pues \(\omega\) tiene orden
nueve. Por ello la representación es fiel. La comparación de los
productos en los dos órdenes da el factor \(bc-ad\).
Un elemento central debe satisfacer \(bc-ad=0\) para todo \((c,d)\).
Tomando sucesivamente \((c,d)=(1,0),(0,1)\), se obtiene
\(b=a=0\). Esto identifica el centro y prueba la exactitud.
\end{proof}

La fase de Wilson para la orientación
\eqref{x_reciprocidad:iv:eq:weyl-area-orientada} es
\(\omega^{\sigma_{\mathrm{ar}}(u,v)}\); el conmutador de la
sección ordenada \(X^aZ^b\) publica su inversa. En particular,
\(c_{\mathrm W}(u,v)-c_{\mathrm W}(v,u)
=-\sigma_{\mathrm{ar}}(u,v)\).
Esta identidad relaciona dos convenciones determinadas, sin identificar
el cociclo polarizado con la forma alternante. La extensión pertenece
al grupo de traslaciones del soporte APP. Su representación actúa sobre
la fibra de fase; el registro de la trayectoria y los acarreos conservan
su residencia en el estado TPK que produjo esa fibra. El orden \(729\)
del grupo y el cardinal \(729\) del ambiente ternario conservan los
tipos distintos establecidos en \eqref{x_reciprocidad:iv:eq:hilbert-finito}.
El centro \(\ZZ/9\ZZ\) de esta extensión registra la fase
del conmutador. La memoria \(C_{12}\) de la sucesión
\(0\to C_{12}\to C_{108}\to C_9\to0\), construida en la
proposición~\ref{x_reciprocidad:prop:k-extension}, pertenece al calendario
dodecafásico. Cada extensión conserva su ley de composición y
su proyección propias.

\begin{proposition}[Álgebra generada por fase y traslación]
\label{x_reciprocidad:iv:prop:weyl-completo}
Los \(81\) operadores \(X^aZ^b\), con \(0\leq a,b<9\), forman
una base de \(\mathcal A_1=M_9(\CC)\). En particular, la
representación de \eqref{x_reciprocidad:iv:eq:weyl-operadores} es irreducible.
\end{proposition}
\begin{proof}
Si \(a\neq0\pmod9\), la matriz \(X^aZ^b\) carece de entradas
diagonales. Si \(a=0\), su traza es
\(\sum_{r=0}^8\omega^{br}\), igual a \(9\) para \(b=0\) y a
\(0\) en los demás casos. Por \eqref{x_reciprocidad:iv:eq:weyl-relacion}, el
producto \((X^aZ^b)^*X^{a'}Z^{b'}\) es un factor de módulo uno
por \(X^{a'-a}Z^{b'-b}\). Por tanto,
\[
 \tau_1\bigl((X^aZ^b)^*X^{a'}Z^{b'}\bigr)
 =\delta_{a,a'}\delta_{b,b'}.
\]
Son \(81\) matrices linealmente independientes en un espacio de
dimensión \(81\). Generan toda el álgebra matricial, cuyo conmutante
está formado por los escalares; de ello resulta la irreducibilidad.
\end{proof}

\begin{theorem}[Unicidad con carácter central primitivo fijado]
\label{x_reciprocidad:iv:thm:weyl-unicidad}
Toda representación unitaria irreducible del grupo presentado por
\(X^9=Z^9=c^9=I\), \(c\) central y \(ZX=cXZ\), en la que
\(c\) actúa como \(\omega I\), es unitariamente equivalente a
\eqref{x_reciprocidad:iv:eq:weyl-operadores}.
\end{theorem}
\begin{proof}
Sea \(v\neq0\) un autovector de \(Z\), de autovalor \(\lambda\).
Como \(Z^9=I\), \(\lambda\) es una potencia de \(\omega\).
Los vectores \(X^kv\), \(0\leq k<9\), tienen autovalores
\(\omega^k\lambda\), distintos porque el carácter es primitivo.
Son ortogonales, y su espacio generado es invariante por \(X\),
\(Z\) y el centro. La irreducibilidad obliga a que sea todo el
espacio. Reordenando la base para comenzar en el autovalor \(1\),
\(Z\) es diagonal con entradas \(1,\omega,\ldots,\omega^8\),
y \(X\) es la traslación cíclica. La normalización de la base
produce la equivalencia unitaria afirmada.
\end{proof}

La primitividad del carácter y la irreducibilidad forman parte del
enunciado. A profundidad \(d\), los operadores que actúan en factores
distintos conmutan, y los productos tensoriales de las bases anteriores
generan \(\mathcal A_d\). La estructura de fase queda así
representada a cada nivel sin identificar el retorno de fase con el
retorno del estado genealógico completo.

\subsection{Dos elevaciones y dos normalizaciones diferentes}

La factorización \eqref{x_reciprocidad:iv:eq:factorizacion-hilbert} permite conservar
toda la fibra nueva o fijar una de sus etiquetas. Para
\(p_j=|\delta_j\rangle\langle\delta_j|\), definimos
\begin{equation}
 \iota_d(a)=a\otimes I_9,
 \qquad \jmath_{d,j}(a)=a\otimes p_j.
 \label{x_reciprocidad:iv:eq:elevaciones}
\end{equation}
\begin{proposition}[Conservación de unidad y traza]
\label{x_reciprocidad:iv:prop:unidad-traza}
La inclusión \(\iota_d\) es unital y conserva la traza normalizada.
La inclusión \(\jmath_{d,j}\) es una inclusión en una esquina y
satisface
\[
 \tau_{d+1}(\iota_d(a))=\tau_d(a),\quad \iota_d(I)=I,
 \qquad
 \tau_{d+1}(\jmath_{d,j}(a))=\frac{\tau_d(a)}9,
 \quad \jmath_{d,j}(I)=I\otimes p_j.
\]
\end{proposition}
\begin{proof}
La traza de un producto tensorial es el producto de trazas.
Como \(\operatorname{Tr}(I_9)=9\) y
\(\operatorname{Tr}(p_j)=1\), dividir por \(9^{d+1}\)
produce las dos fórmulas. Las imágenes de la identidad se leen
directamente en \eqref{x_reciprocidad:iv:eq:elevaciones}.
\end{proof}

La inclusión unital conserva el conjunto de las \(9\) fases y su
normalización global. La esquina conserva una continuación marcada y
su proyector de soporte. Ambas operaciones son legítimas, pero
transmiten datos distintos a la siguiente profundidad.

\subsection{Cierre inductivo y realización tracial}

\begin{theorem}[Álgebra uniformemente hiperfinita nonádica]
\label{x_reciprocidad:iv:thm:uhf}
El cierre normado de la unión mediante las inclusiones \(\iota_d\)
es
\[
 \mathcal A_{9^\infty}
 =\varinjlim(\mathcal A_d,\iota_d)
 \cong\bigotimes_{n=1}^{\infty}M_9(\CC).
\]
Es un álgebra UHF de tipo \(9^\infty=3^\infty\), con una única
traza normalizada \(\tau\), cuya restricción a \(\mathcal A_d\)
es \(\tau_d\).
\end{theorem}
\begin{proof}
Las inclusiones son inyectivas y unitales. La clausura de la unión
creciente de álgebras matriciales es precisamente la construcción UHF.
La compatibilidad de la proposición \ref{x_reciprocidad:iv:prop:unidad-traza}
define \(\tau\) sobre la unión; positividad y norma uno permiten
extenderla al cierre. Toda traza normalizada del límite se restringe
a la única traza normalizada de cada álgebra matricial. Como la unión
es densa, esas restricciones determinan la traza de manera única.
\end{proof}

Sea \(\mathcal H_\tau\) la completación de
\(\mathcal A_{9^\infty}\) para
\(\langle a,b\rangle_\tau=\tau(a^*b)\), y sea \(\pi_\tau\)
su representación por multiplicación izquierda. El vector definido
por la identidad se denota \(\Omega_\tau\).

\begin{theorem}[Factor tracial del refinamiento completo]
\label{x_reciprocidad:iv:thm:factor-tracial}
La clausura débil
\(\mathcal R_9=\pi_\tau(\mathcal A_{9^\infty})''\)
es un factor hiperfinito de tipo \(II_1\). Las inclusiones
nonádicas conservan en él la unidad y la traza normalizada.
\end{theorem}
\begin{proof}
La traza producto se extiende a una traza normal fiel y finita
en la representación tracial. Para ver que el centro es trivial,
sea \(E_d\) la expectativa tracial sobre \(\pi_\tau(\mathcal A_d)\).
En un tensor de longitud \(n>d\), toma la traza de los últimos
\(n-d\) factores. Estas aplicaciones se prolongan por continuidad
a las proyecciones ortogonales correspondientes en \(L^2\).
Si \(z\) es central, \(E_d(z)\) conmuta con
\(\pi_\tau(\mathcal A_d)\), por la bimodularidad de \(E_d\).
Es por ello escalar e igual a \(\tau(z)I\). La densidad de la
unión matricial implica \(E_d(z)\to z\) en \(L^2\), de donde
\(z=\tau(z)I\). El álgebra es, pues, un factor finito.
Contiene unitariamente matrices de tamaños \(9^d\) sin cota,
por lo que no es un factor matricial finito. Es de tipo \(II_1\).
Es hiperfinita porque esa misma unión de matrices es débilmente
densa. La conservación de unidad y traza procede de cada inclusión
finita y se mantiene al tomar las clausuras.
\end{proof}

\subsection{Rangos enteros y dimensión tracial continua}

En \(\mathcal A_d\), una proyección \(P\) tiene dimensión
normalizada \(\tau_d(P)=\operatorname{rank}(P)/9^d\).
Los rangos siguen siendo enteros a toda profundidad finita. Sus
normalizaciones compatibles permiten el siguiente paso al límite.

\begin{proposition}[Realización de cada dimensión tracial]
\label{x_reciprocidad:iv:prop:dimension-tracial}
Para cada \(t\in[0,1]\) existe una proyección
\(P_t\in\mathcal R_9\) con \(\tau(P_t)=t\).
\end{proposition}
\begin{proof}
Sea \(m_d=\lfloor9^dt\rfloor\). Se tiene
\(0\leq m_{d+1}-9m_d\leq8\). Elegida una proyección diagonal
de rango \(m_d\), su inclusión tiene rango \(9m_d\); se
añaden \(m_{d+1}-9m_d\) posiciones diagonales en su complemento.
Se obtiene así una sucesión creciente de proyecciones dentro de
\(\mathcal R_9\). Su límite fuerte es una proyección \(P_t\),
y la normalidad de la traza da
\[
 \tau(P_t)=\lim_{d\to\infty}\frac{\lfloor9^dt\rfloor}{9^d}=t.
\]
\end{proof}

Aquí \(t\) designa la coordenada de la dimensión que se representa
en el cierre, no una entrada del generador APP. La afirmación
describe la amplitud de la realización tracial resultante. Para
compararla con un observable determinado sigue siendo necesario
identificar la proyección y la información genealógica que conserva.

Por contraste, las isometrías
\(V_{d,j}\xi=\xi\otimes\delta_j\) realizan
\(\jmath_{d,j}(a)=V_{d,j}aV_{d,j}^*\). En el límite de una
ruta marcada, los espacios finitos son densos. Las esquinas
aproximan todos los operadores de rango finito; su cierre normado
es \(\mathcal K(\mathcal H_\infty)\), y su bicommutante es
\(B(\mathcal H_\infty)\), de tipo \(I_\infty\). Esta
diferencia conserva el contenido operatorio de las dos elecciones
de \eqref{x_reciprocidad:iv:eq:elevaciones}: la selección de una continuación y
la prolongación unital de toda la fibra tienen límites diferentes.

La realización hilbertiana, sus operadores de fase y el cierre
tracial quedan construidos en esta sección, con sus elecciones de
inclusión explícitas. Constituyen el
soporte operatorio que se utilizará después para el intercambio de
coordenadas y para el Hamiltoniano compacto.
